% Created 2026-06-19 Fri 01:47
% Intended LaTeX compiler: pdflatex
\documentclass{article}
\usepackage[utf8]{inputenc}
\usepackage[T1]{fontenc}
\usepackage{graphicx}
\usepackage{longtable}
\usepackage{wrapfig}
\usepackage{rotating}
\usepackage[normalem]{ulem}
\usepackage{amsmath}
\usepackage{amssymb}
\usepackage{capt-of}
\usepackage{hyperref}

\usepackage{fancyhdr}
\usepackage[top=25mm, left=25mm, right=25mm]{geometry}
\usepackage{longtable}
\fancyhead[R]{}
\setlength\headheight{43.0pt}
\usepackage[T1]{fontenc}
\usepackage[utf8]{inputenc}
\usepackage[spanish, ]{babel}
\usepackage[backend=biber]{biblatex}
\addbibresource{/home/josetm-epn/MI-REPO/EPN-Lectures/iccd332ArqComp-2024-B/FormatoTareas/bibliography.bib}
\author{Joset Nicoly Muela Criollo}
\date{2026-06-17}
\title{Plantilla para Toma de Notas en Clase}
\hypersetup{
 pdfauthor={Joset Nicoly Muela Criollo},
 pdftitle={Plantilla para Toma de Notas en Clase},
 pdfkeywords={},
 pdfsubject={},
 pdfcreator={Emacs 27.1 (Org mode 9.7.5)}, 
 pdflang={Español}}
\begin{document}

\maketitle
\tableofcontents

\fancyhead[C]{\includegraphics[scale=0.05]{./logoEPN.jpg}\\
ESCUELA POLITÉCNICA NACIONAL\\FACULTAD DE INGENIERÍA DE SISTEMAS\\
ARQUITECTURA DE COMPUTADORES}
\thispagestyle{fancy}
\section{2026-06-17 Elemento de Diseño de un Bus}
\label{sec:org8800bb0}
\subsection{Preguntas y Términos clave (Cue):}
\label{sec:org000433e}
\begin{description}
\item[{¿Qué clasifica a un bus dedicado de uno multiplexado?}] 

\item[{Tipos de Arbitraje}] 

\item[{Temporización Síncrona vs Asíncrona}] 

\item[{Importancia del ancho del bus}] 

\item[{Tipos de transferencia de datos}] 
\end{description}
\subsection{Notas (Notes):}
\label{sec:org2a77723}
\begin{itemize}
\item \textbf{Clasificación por Tipo de Líneas}:
\begin{itemize}
\item \emph{Dedicadas}: Líneas físicas asignadas perman
entemente a una función ,buses independientes para datos y
direcciones. Ofrecen alto rendimiento pero incrementan el espacio
y costo.
\item \emph{Multiplexadas en el tiempo}: Las mismas líneas físicas se
comparten para distintas funciones en instantes de tiempo
diferentes, la misma línea envía la dirección al inicio del
ciclo y luego el dato. Ahorra pines y espacio, pero ralentiza el
sistema al requerir circuitos lógicos más complejos.
\end{itemize}
\item \textbf{Métodos de Arbitraje (Control del medio compartido)}:
\begin{itemize}
\item \emph{Centralizado}: Existe un único hardware dedicado, llamado
controlador del bus o árbitro, que recibe las peticiones y concede
el acceso al bus.
\item \emph{Distribuido}: No existe un nodo central. Los módulos incluyen
lógica de control propia y cooperan conjuntamente para decidir el
acceso mediante algoritmos distribuidos.
\end{itemize}
\item \textbf{Esquemas de Temporización}:
\begin{itemize}
\item \emph{Síncrona}: Las acciones en el bus están regidas por una señal de
reloj global (tren de pulsos lógicos \(0\) y \(1\)). Todos los
dispositivos se sincronizan con este reloj. Es rápida y simple,
pero obliga a que todo el bus trabaje a la velocidad del
dispositivo más lento.
\item \emph{Asíncrona}: No depende de un reloj. La secuencia de eventos
ocurre mediante un protocolo de "saludo" o \emph{handshaking} (un
evento dispara el siguiente: Petición \(\rightarrow\)
Confirmación). Es altamente flexible para mezclar dispositivos
rápidos y lentos, pero su hardware de control es complejo.
\end{itemize}
\item \textbf{Ancho del Bus}:
\begin{itemize}
\item \emph{Bus de Datos}: Un mayor número de líneas físicas paralelas
permite transmitir palabras de datos más grandes en un solo ciclo,
impactando directamente en el rendimiento global.
\item \emph{Bus de Direcciones}: El ancho de este bus limita físicamente el
rango máximo de memoria principal direccionable directamente por
la CPU (\(2^{\text{ancho}}\) bytes).
\begin{itemize}
\item \textbf{Tipos de Transferencia de Datos}:
\end{itemize}
\item \emph{Lectura} y \emph{Escritura} clásicas de memoria o E/S.
\item \emph{Lectura-Modificación-Escritura}: Operación atómica (no
interrumpible) donde se lee un dato, se modifica y se reescribe
inmediatamente; fundamental para la sincronización de hilos o
semáforos en el Sistema Operativo.
\end{itemize}
\end{itemize}

\subsection{Enlaces de Interés}
\label{sec:org0cc9ab5}
\begin{itemize}
\item\relax [[\url{https://lsc.cornell.edu/how-to-study/taking-notes/cornell-note-taking-system/}][
\item \href{https://biblioteca.univalle.edu.ni/files/original/c1b1f5d1c12abc60a246e2a0d784f7c9d163ee81.pdf}{Libro William Stallings}
\end{itemize}
\subsection{Resumen (Summary):}
\label{sec:org7229102}
El diseño de un bus implica un balance de ingeniería entre costo,
espacio físico y rendimiento computacional. Los parámetros clave
definidos por Stallings determinan la flexibilidad y velocidad del
sistema: el arbitraje gestiona la competencia por el recurso
compartido; la temporización define el método de coordinación lógica
(reloj rígido vs handshake flexible); y el tipo de líneas (dedicadas o
multiplexadas) junto al ancho físico determinan directamente las
capacidades máximas de transferencia de datos y de direccionamiento de
memoria del computador.



\section{{\bfseries\sffamily TODO} Mis Tareas [100\%]}
\label{sec:orgbdd17e2}
\begin{itemize}
\item[{$\boxtimes$}] Leer este documento
\item[{$\boxtimes$}] Escribir el resumen  de la exposición 1
\end{itemize}

\printbibliography
\end{document}
